% Generated draft: requires mathematical and source review.
\documentclass[11pt]{article}
\usepackage[margin=25mm]{geometry}
\usepackage{fontspec}
\setmainfont{Latin Modern Roman}
\usepackage{amsmath,amssymb,mathtools,unicode-math}
\setmathfont{Latin Modern Math}
\usepackage{enumitem,needspace,xurl,hyperref}
\setlength{\parindent}{0pt}
\setlength{\parskip}{5pt}
\setlength{\emergencystretch}{3em}
\newcommand{\N}{\mathbb{N}}
\newcommand{\Z}{\mathbb{Z}}
\newcommand{\Q}{\mathbb{Q}}
\newcommand{\R}{\mathbb{R}}
\newcommand{\C}{\mathbb{C}}
\newcommand{\F}{\mathbb{F}}
\newcommand{\A}{\mathbb{A}}
\newcommand{\PP}{\mathbb{P}}
\newcommand{\E}{\mathbb{E}}
\newcommand{\eps}{\varepsilon}
\newcommand{\dd}{\,\mathrm{d}}
\begin{document}
\section*{UnsolvedMath ID: 30001044}
\textbf{Source problem code:} catalogue entry
\subsection{Definitions and mathematical statement}
Let \(M\) be a compact manifold of the source's unspecified ``hyperbolic type,'' equipped with an arbitrary Riemannian metric \(g\), and let \(SM\) be its unit tangent bundle. The set \(SM_{\min}^g\) is described as the velocities of globally \(g\)-minimizing geodesics. A necessary convention to verify is whether minimizing means that a lift \(\widetilde\gamma\) to the universal cover satisfies
\[
d_{\widetilde g}(\widetilde\gamma(s),\widetilde\gamma(t))=|s-t|\quad\text{for every }s,t\in\R.
\]
This is distinct from global minimization in the compact manifold itself and is relevant to the possibility of closed geodesics.

The source asks:
\begin{enumerate}
\item Does \(SM_{\min}^g\) contain the velocity orbit of a closed hyperbolic geodesic? Hyperbolicity here ordinarily means that the transverse return map has no eigenvalue on the unit circle, not that \(g\) has negative curvature.
\item Is there an ergodic geodesic-flow-invariant probability measure with support contained in \(SM_{\min}^g\), hyperbolic in the sense of having no zero transverse Lyapunov exponents? The unavoidable exponent in the flow direction is excluded in that convention.
\item Separately, for an area-preserving surface diffeomorphism \(f\), can its intended measure-preserving system be measurably conjugated to \(R_\beta:S^1\to S^1\), \(R_\beta(x)=x+\beta\pmod1\), where \(\beta\) is Diophantine? A usual meaning of Diophantine is that \(|\beta-p/q|\geq c q^{-\tau}\) for some \(c>0\), finite \(\tau\), and every \(p\in\Z,q\geq1\). Measurable conjugacy is an invertible measure-preserving intertwining map modulo null sets.
\end{enumerate}
The last clause supplies neither a restriction on \(f\) beyond area preservation nor the invariant probability measure, domain subset, or relation to the geodesic questions. Its quantifier is not precise enough to assert such a rotation model for every area-preserving diffeomorphism. No missing dynamics hypotheses are invented.
\subsection{Short English statement}
Ask for hyperbolic periodic orbits and hyperbolic invariant measures among minimizing geodesics, and retain a separate question about measurable Diophantine rotation models. The minimizing-lift convention, hyperbolic-type assumption and hypotheses of the surface-map clause need clarification.
\subsection{Sources}
\begin{itemize}
\item Local UnsolvedMath registry, permanent ID \(30001044\), source code \texttt{catalogue entry}, ``Hyperbolic Minimal Geodesics and Measurable Rotation Models.'' Registry curation: Ulam AI / UnsolvedMath Contributors, CC BY 4.0; underlying sources retain their supplied provenance.
\item \emph{Geometric Group Theory, Hyperbolic Dynamics and Symplectic Geometry} (2009), Oberwolfach Reports, pp.\ 2337--2406, DOI \url{https://doi.org/10.4171/owr/2008/42}.
\item Supplied literature triage, \texttt{literature\_checked\_at}: 2026-08-21: labeled open; says the extracted questions concerning closed hyperbolic geodesics, invariant measures, and measurable rotation models on minimal geodesics remain unresolved and no authoritative resolution was found. This is supplied dated curation, not independently verified, and does not fill the missing surface-map hypotheses.
\end{itemize}
\end{document}
